\paragraph{A fixed-radius inverse for every current-order block.}
The preceding Neumann argument used $\|K_n\|\le C_K/(n+1)$ and therefore treated finitely many small orders separately. That step is unnecessary if one retains radial valuation. Each term of the current-order operator $K_n$ has a final radial inverse $J_1$ or $J_2$; the pressure-feedback channel additionally has an axis primitive. The similarity order $n$ remains fixed while $K_n$ is iterated, but the radial valuation of the input increases.

\begin{theorem}[Factorially convergent current-order resolvent]\label{thm:factorial-resolvent}
Fix a common analytic coefficient space $\mathcal B_{R,\rho}$, and assume the finite leading-profile multiplier norms of Theorem~\ref{thm:fourblock}. Let $K_n$ denote the exact two-by-two operator (47)--(50). A finite constant $C_*$, independent of $n$ and of the input valuation $m$, exists such that, for $m\ge n\ge1$,
\begin{equation}\label{eq:valuation-tame-K}
 K_n(I_m\oplus I_m)\subseteq I_{m+1}\oplus I_{m+1},
 \qquad\|K_nF\|\le\frac{C_*}{m+1}\|F\|.
\end{equation}
Consequently,
\begin{equation}\label{eq:factorial-K}
\|K_n^\ell\|_{I_n^2\to I_{n+\ell}^2}
\le\frac{C_*^\ell n!}{(n+\ell)!}
\le\frac{C_*^\ell}{\ell!},
\end{equation}
and, for every $n\ge1$, $I-K_n$ is invertible on $I_n^2$ with
\begin{equation}\label{eq:fixed-radius-inverse}
 (I-K_n)^{-1}=\sum_{\ell\ge0}K_n^\ell,
 \qquad\|(I-K_n)^{-1}\|\le e^{C_*}.
\end{equation}
There is no radius-dependent exceptional startup order.
\end{theorem}
\begin{proof}
The source identities \eqref{eq:swirlblock}, \eqref{eq:Voperator}, \eqref{eq:phiu}, \eqref{eq:uu}, and \eqref{eq:pressureblock} express each block as a finite sum of complete radial-inverse composites. The estimates \eqref{eq:comp0}--\eqref{eq:compX} hold on arbitrary $I_m$: an undifferentiated term gains $(m+1)^{-2}$ and one $D_X$ or $\partial_\eta$ costs at most one power of $m+1$. The similarity coefficients cause no loss uniform in $m\ge n$, since
\[
 |b+\lambda_n|\le(|b|+2h)(m+1),\quad
 |c+\lambda_n|\le(|c|+2h)(m+1),\quad
 |D+\lambda_n|\le(|D|+2h)(m+1).
\]
The radial-velocity cross blocks include an average and satisfy
$\|\mathcal V_{n,\nu}^{g,q}|_{I_m}\|\le C_V(g,q)/(m+1)^2$.
The pressure block includes an axis primitive, followed by $J_1$;
using \eqref{eq:averageprim} at input valuation $m$ gives
$\|K_n^{U\varphi}|_{I_m}\|\le C_{U\varphi}/(m+1)^2$.
Fixed leading-profile multipliers have nonnegative radial degrees;
the final inverses increase valuation by at least one.
Thus $C_*=1+\max\{C_{\varphi\varphi}+C_{U\varphi},
 C_{UU}+C_V(m_\varphi,q_\varphi)\}$, with the constants
defined in Section~4, is sufficient.

Successively apply the gain at valuations $n,n+1,\ldots,n+\ell-1$.
The factorial inequality in \eqref{eq:factorial-K} follows because
$(n+1)\cdots(n+\ell)\ge\ell!$. The Neumann series converges
absolutely in operator norm, without the assumption $\|K_n\|<1$.
Its partial sums are inverses modulo the remainder $K_n^{N+1}$,
which tends to zero. Summation gives the bound $e^{C_*}$.
\end{proof}

\paragraph{A two-shift convolution majorant.}
The pressure source contains positive-order quadratic sums at
$i+j=n$ and $i+j=n-1$; these cannot be merged without an assumption
on the coefficient sequence.

\begin{proposition}[Constructive majorant with both quadratic shifts]
\label{prop:two-shift-majorant}
Suppose nonnegative $B_n$ are finite at $n=1,2$ and satisfy, for $n\ge3$,
\begin{equation}\label{eq:two-shift-recurrence}
 B_n\le H\left(
 B_{n-1}+B_{n-2}+
 \sum_{\substack{i+j=n\\i,j\ge1}}B_iB_j+
 \sum_{\substack{i+j=n-1\\i,j\ge1}}B_iB_j\right),
\end{equation}
with fixed $H\ge0$. Put $M=[32(1+H)]^{-1}$ and choose
\[
 T\ge\max\left\{2,16H,6\sqrt H,\frac{4B_1}{M},
 3\sqrt{\frac{B_2}{M}}\right\}.
\]
Then $B_n\le MT^n/(n+1)^2$ for all $n\ge1$.
\end{proposition}
\begin{proof}
The startup bounds follow from the choice of $T$.
Apply the inductive hypothesis and \eqref{eq:convolution}.
For $n\ge3$ the right side of \eqref{eq:two-shift-recurrence},
divided by $MT^n/(n+1)^2$, is at most
\[
 \frac{2H}{T}+\frac{4H}{T^2}
 +4HM\left(1+\frac1T\right)<1.
\]
The $n-1$ convolution contributes the final term $4HM/T$;
it is not identified with the $n$ convolution. Induction closes.
\end{proof}

\paragraph{Source-to-majorant reduction on a prescribed strip.}
Put $B_n=\|\varphi_n\|_{\mathcal B}+\|U_n\|_{\mathcal B}$.
Let $C_{\rm del,mom}$ and $C_{\rm mom}$ denote the direct delayed
and bilinear momentum constants above, $C_E$ the shifted
pressure-to-velocity bound, and
$C_L=C_{\rm P,time}+C_{\rm P,lap}+C_{\rm lead,P}$,
$C_D=C_{\rm P,del}$, $C_{\rm sw}$ and $C_{\rm quad,P}$
the shifted-pressure source constants in
\eqref{eq:complete-two-space-P-ledger}. After separating
$\varphi_0\varphi_n$ into the current-order block $K_n^{U\varphi}$,
the assembled lower-order estimates have the form
\begin{equation}\label{eq:fixed-radius-source-budget}
 \|S_nf_n\|\le L_1B_{n-1}+L_2B_{n-2}
 +Q_0\!\sum_{\substack{i+j=n\\i,j\ge1}}B_iB_j
 +Q_1\!\sum_{\substack{i+j=n-1\\i,j\ge1}}B_iB_j,
\end{equation}
with conservative budgets
\[
 L_1=2C_{\rm del,mom}+C_EC_L,\quad
 L_2=C_EC_D,\quad
 Q_0=2C_{\rm mom}+C_EC_{\rm sw},\quad
 Q_1=C_EC_{\rm quad,P}.
\]
They are independent of $n$ on one preselected common analytic
space. The source decomposition
\eqref{eq:lower-swirl}--\eqref{eq:lower-pressure}
and the previous component theorems provide this budget only
upon verifying their exact normalization and leading-profile
hypotheses for the proposed unforced hierarchy.

Under these stated hypotheses the inverse theorem gives
\eqref{eq:two-shift-recurrence} with
\begin{equation}\label{eq:source-majorant-H}
 H=e^{C_*}\max\{1,L_1,L_2,Q_0,Q_1\}.
\end{equation}
Proposition~\ref{prop:two-shift-majorant} then gives geometric
velocity coefficients on $T|q^{2h}|<2$ and $|X|<R/8$.
The shifted pressure ledger implies
$\|\Pi_n\|_{\mathcal P}\le\Gamma T^n/(n+1)^2$
with finite $\Gamma$ as in \eqref{eq:conditional-pressure-geometric}.
The physical coordinate Jacobian is $L(\eta)>0$ on the compact
subcore $q>0$, so each fixed physical derivative passes through
the normally convergent series once the exact startup source
equations and chain rule have been checked.

\paragraph{Scope of this result.}
Theorem~\ref{thm:factorial-resolvent} removes the circular
finite-startup radius cutoff, provided the analytic strip and all
leading multipliers have been fixed beforehand.
Proposition~B.2 of \cite{openai2026} gives an analytic source
leading profile for admissible parameters. The identification
of that particular profile, its pressure datum, and the unforced
hierarchy's startup normalization require independent correspondence
checks. These results are a conditional local reconstruction criterion,
not a proof of the proposed unforced Theorem~1.1 or a global singularity.
